% Block 3 of the slides (03_context.tex).
\section{Context, effort and persistent instructions}\label{sec:context}

\onslides{the context window; context rot; context hygiene; the cost of thinking;
instruction files; skills; a prompt as a specification; the workflow explore--plan--%
implement--verify--commit. This block is the first to shorten if time runs out.}

\subsection{The context window}

The context is the token sequence $x_1, \ldots, x_n$ of Section~\ref{sec:llm}. In an agent
session it contains, in order: the harness's system instructions and tool descriptions,
your instruction files, your messages, every file the agent read, every command output,
and the model's own answers and reasoning. Its length is capped by the context window,
currently $10^5$ to $10^6$ tokens (the largest Llama~3 model has $1.28 \times 10^5$
\rf{llama3}). At about $0.7$ words per token, $10^5$ tokens is a long novel. That sounds
like a lot until an agent reads a few source files, a test log and a long solver output.

\subsection{Context rot}

Accuracy degrades as the context grows, well before the window is full. Hong, Troynikov
and Huber \rf{contextrot} tested 18 models on tasks that stay equally hard as the input
grows, so that any drop is due to length alone. Every model degraded with input length,
and non-uniformly. Three findings are worth quoting: a single distractor (text similar to
the answer but wrong) lowered accuracy, and its effect grew with length; on a
conversational memory benchmark, all models did much better with a focused input of about
300 tokens than with the full $\sim 113$K-token history; and models performed
\emph{worse} when the irrelevant filler was a coherent text than when its sentences were
shuffled. Earlier, Liu et al.\ \rf{lostmiddle} found that models use information at the
beginning and end of a long context better than information in the middle.

Practical consequences:
\begin{itemize}
  \item A wrong turn stays in the context and keeps influencing later answers. Arguing
        with the agent adds more of it. Start a fresh session with a better prompt.
  \item When the window fills, harnesses \emph{compact}: they replace the history by a
        summary written by the model \rf{context-eng}. The summary is lossy; details you
        care about may not survive.
  \item One task, one session. Write what must survive into files (below).
\end{itemize}

\textbf{Point, do not paste.} Three thousand lines of solver log in your message cost
context and bury the one relevant line. Naming the failing test and asking the agent to
run only that test lets it read what it needs. Anthropic's engineering team calls this
loading context ``just in time'' through references such as file paths \rf{context-eng}.

\subsection{What a session costs}

Providers charge per token, at different prices for input and output. With $N_{\mathrm{in}}$
and $N_{\mathrm{out}}$ the numbers of input and output tokens and $c_{\mathrm{in}}$,
$c_{\mathrm{out}}$ the prices per token,
\begin{equation}
  \text{cost} \approx N_{\mathrm{in}}\, c_{\mathrm{in}} + N_{\mathrm{out}}\, c_{\mathrm{out}}.
\end{equation}
The slide's remark that every agent step re-reads the context has a quantitative
consequence. Suppose a session takes $K$ steps (tool calls), starts from a context of
$C_0$ tokens, and each step adds about $a$ tokens (a file read, a command output). Step
$k$ reads the whole context, $C_0 + k a$ tokens, so
\begin{equation}
  N_{\mathrm{in}} \approx \sum_{k=1}^{K} (C_0 + k a) \approx K C_0 + \tfrac12 a K^2 ,
\end{equation}
quadratic in the length of the session. Long sessions are expensive as well as less
accurate. Many providers discount input that repeats an earlier prefix (\emph{prompt
caching}), which softens but does not remove the effect.

$N_{\mathrm{out}}$ includes the reasoning tokens of Section~\ref{sec:llm}, which you may
never see. The \emph{effort} setting of current agents controls how much the model
reasons before answering: low effort for mechanical edits, high for debugging and
derivations. For scale, the TheoPhysEval agents solved day-long research problems for
under 3\,USD each \rf{theophyseval}.

\subsection{Instruction files}

Because the model remembers nothing between sessions, harnesses read an instruction file
from the project folder at the start of every session and put it into the context. The
file name depends on the tool: \cmd{AGENTS.md} is an open convention read by many agents
\rf{agentsmd}; Claude Code reads \cmd{CLAUDE.md} \rf{cc-best}. The \cmd{DrivenSpin/}
project has an \cmd{AGENTS.md} and a one-line \cmd{CLAUDE.md} that points to it.

What belongs in it: the conventions the agent cannot guess (units, basis order, sign of
$\Delta$), the command that runs the tests, and the rules that protect you (never weaken a
test, never write to \cmd{data/exp\_raw}). What does not: anything the agent can read from
the code, long tutorials, information that changes often. The file is in every context,
so every line costs tokens in every step; a long file also dilutes the rules that matter
\rf{cc-best}.

\begin{placeholder}[Placeholder: \texttt{DrivenSpin/AGENTS.md}]
Quote the real file here once D3 exists. The slide shows an illustrative excerpt:
conventions ($\hbar = 1$, \textmu s, rad/\textmu s, basis $(\ket{e}, \ket{g})$,
$H = (\Delta/2)\sigma_z + (\Omega/2)\sigma_x$), the test command
\verb|julia --project -e 'using Pkg; Pkg.test()'|, ``tests first'', ``never modify
\texttt{test/} while implementing'', ``never write to \texttt{data/exp\_raw}''.
\end{placeholder}

\subsection{Skills}

A \emph{skill} is a folder with a file \cmd{SKILL.md} that starts with a name and a
one-line description, followed by instructions, and optionally scripts and further
files. At the start of a session only the names and descriptions enter the context; when a
task matches a description, the agent reads the full file, and it opens further files
only when a step needs them \rf{skills}. This \emph{progressive disclosure} lets a lab keep
many procedures (fit Ramsey data, calibrate a pulse, format a figure) at the cost of one
line each. The slide's \cmd{fit-ramsey} skill is an example of the form, not a file in the
repository. Instruction files and skills are \emph{procedural memory written to disk}:
the model does not learn, but the folder does.

\subsection{A prompt is a specification}

The slide contrasts ``Simulate a qubit with decoherence'' with a prompt that names the
goal and context (which file, which conventions), the verification (the tests that define
``done''), and the constraints (what must not change; plan before code). Each part closes
a door the agent would otherwise walk through: without the file it invents a new one;
without the conventions it picks its own; without the tests ``done'' means ``looks
done''; without constraints it ``improves'' the Hamiltonian.

\subsection{A workflow that works}

Explore (read the code), plan (you approve), implement (small steps), verify (tests,
plots), commit (you read the diff). The same sequence is the core recommendation of
Anthropic's guide \rf{cc-best}, and it is tool-independent. Most failures come from
skipping \emph{plan} (the agent solves the wrong problem) or \emph{verify} (it solves the
problem wrongly).

\begin{qa}
\Q{How big is my context right now?}
Most agents can show it (a command or a status line). If you cannot see it, start fresh
sessions often.

\Q{Should I always use the highest effort?}
No: it costs more and is slower, and for mechanical edits it does not help. Raise it when
the agent goes in circles on a hard problem.

\Q{Instruction file or skill?}
Rules that apply to every task go into the instruction file. Procedures that apply to
some tasks go into skills, so they cost context only when used.

\Q{Is ``context rot'' a real effect or a vendor's term?}
The term comes from the Chroma report \rf{contextrot}, a controlled measurement on 18
models; the effect is consistent with earlier academic work \rf{lostmiddle}.
\end{qa}

\textbf{Read next:} \emph{Effective context engineering for AI agents} \rf{context-eng};
the Chroma report \rf{contextrot}, whose figures are worth a look; the skills article
\rf{skills}.
